% \subsection{Local systems and the Reidemeister trace}

% Let $g\colon \anima\to \anima$ be a self-map of a compact anima $\anima$. Consider again the lax commutative square~\eqref{eq:lax-comm-square-lefschetz} from~\cref{sec: lefschetz theory}, which we viewed as a morphism in the category of traceable endomorphisms $\Endtrl(\Prlstf)$. If $g\colon \anima\to \anima$ happened to be induced by a self-map $f\colon\tset\to\tset$ for some well-behaved space $\tset$, we recovered the Lefschetz fixed point theorem for $f$ by factoring~\eqref{eq:lax-comm-square-lefschetz} through the object $(\D{\tset },f_!)\in\Endtrl(\Prlstf)$.

% Putting topology aside for now, a more obvious factorization of~\eqref{eq:lax-comm-square-lefschetz} is through the category of local systems (aka parametrized spectra) $\Sp^\anima$:
% \begin{prop}
%   Let $g\colon \anima\to \anima$ be a self-map of a compact anima $\anima$. Considered as a morphism in $\Endtrl(\Prlstf)$, the lax commutative square~\eqref{eq:lax-comm-square-lefschetz} factors as
%   \begin{equation}\label{eq:lax-composition-reidemeister}
% \begin{tikzcd}
% 	\Sp & {\Sp^{\anima}} & \Sp \\
% 	\Sp & {\Sp^{\anima}} & \Sp
% 	\arrow["{\anima^{\ast}}", from=1-1, to=1-2]
% 	\arrow["\id"', from=1-1, to=2-1]
% 	\arrow["{\anima_\sharp}", from=1-2, to=1-3]
% 	\arrow["{g^{\ast}}", from=1-2, to=2-2]
% 	\arrow["\id", from=1-3, to=2-3]
% 	\arrow[Rightarrow, from=2-2, to=1-3]
% 	\arrow["{\anima^{\ast}}"', from=2-1, to=2-2]
%     \arrow[Rightarrow, from=2-1, to=1-2]
% 	\arrow["{\anima_{\sharp}}"', from=2-2, to=2-3]
% \end{tikzcd}
%     % \left(
%     %   \begin{tikzcd}[column sep=large]
%     %     \Sp^X \arrow[r,"X_\sharp"] \arrow[d,"g^*"] & \Sp \arrow[d,equal]\\
%     %     \Sp^X \arrow[r,swap,"X_\sharp"]  \arrow[ru,Rightarrow, shorten <=5pt,shorten >=5pt]  & \Sp
%     %   \end{tikzcd}
%     % \right)
%     % \circ
%     % \left(
%     %   \begin{tikzcd}[column sep=large]
%     %     \Sp \arrow[r,"X^*"] \arrow[d,equal] & \Sp^X \arrow[d,"g^*"]\\
%     %     \Sp \arrow[r,swap,"X^*"] & \Sp^X
%     %   \end{tikzcd}
%     % \right),
%   \end{equation}
%   where the left hand $2$-cell is the invertible one coming from functoriality of pullback, and the right hand $2$-cell is the counit transformation $\anima_\sharp g^*\simeq \anima_\sharp g_\sharp g^*\to \anima_\sharp$.
% \end{prop}
% \begin{proof}
%   Aside from the usual adjunctions $\anima_\sharp\dashv \anima^*\dashv \anima_*$, the fact that $\anima$ is compact furnishes an additional adjunction $\anima_*\dashv \Hom_{\Sp^\anima} (D_\anima,\anima^*({-}))$, where $D_\anima$ is the dualizing object of $\anima$. It follows that the horizontal morphisms in both squares are strongly continuous, and hence that they define morphisms in $\Endtrl(\Prlstf)$. Concatenating the two $2$-cells recovers the one in~\eqref{eq:lax-comm-square-lefschetz} on the nose.
% \end{proof}
% As in the previous section, we then get
% \begin{cor}\label{cor: reidemeister lefschetz triangle}
%   For a self-map $g\colon \anima\to \anima$ of a compact anima $\anima$, there is a commutative diagram
%   \begin{equation}
%     \label{eq:reidemeister-lefschetz-triangle}
%     \begin{tikzcd}[column sep=tiny]
%       & \tr{g^*\curvearrowright\Sp^\anima}{}\arrow[dr] \\
%       \SS\arrow[rr]\arrow[ru] && \SS,
%     \end{tikzcd}
%   \end{equation}
%   where the bottom arrow is the Lefschetz trace of $g$.
% \end{cor}
% Just as in the previous section, the axioms of six-functor formalisms allow us to quickly compute the trace $\tr{g^*\curvearrowright\Sp^\anima}{}$ with the suspension spectrum of a space of ``fixed points'' for the map $g$. Here fixed points must be interpreted in a homotopy-invariant way; we let $\cL^g\anima$ denote the homotopical fixed point space defined to be the homotopy pullback
% \[
%   \begin{tikzcd}
%     \cL^{g}\anima  & \anima \\
%     \anima & \anima \times \anima\,.
%     \arrow[from=1-1, to=1-2]
%     \arrow["\lrcorner",phantom,very near start, from=1-1, to=2-2]
%     \arrow[from=1-1, to=2-1]
%     \arrow["\Delta", from=1-2, to=2-2]
%     \arrow["\Gamma^g"', from=2-1, to=2-2]
%   \end{tikzcd}
% \]
% That is, we have
% \begin{lem}\label{lem: reidemeister trace computation}
%   Let $g\colon \anima\to \anima$ be a self-map of a compact anima. There is an equivalence of spectra
%   \[
%     \tr{g^*\curvearrowright\Sp^\anima}{}\simeq\SS\lbrack\cL^{g}\anima\rbrack,
%   \]
%   natural in intertwining maps $h\colon (\anima,g)\to (\anima',g')$. 
% \end{lem}
% \begin{proof}
%   The identification comes about exactly as in the proof of~\cref{lem: trace computation}. Unlike in the topological setting, we do not need to require anything of $h$ beyond being an intertwining map, since any intertwining map induces a morphism
%   \begin{equation}
%     \label{eq:square-induced-by-proper-morphism}
%     \begin{tikzcd}
%       \Sp^{\anima}\arrow[r,"h_\sharp"]\arrow[d,"g^*"]
%       & \Sp^{\anima' }\arrow[d,"(g')^*"]\\
%       \Sp^{\anima}\arrow[r,"h_\sharp"]\arrow[ru,Rightarrow, shorten <=5pt,shorten >=5pt]
%       & \Sp^{\anima'}
%     \end{tikzcd}
%   \end{equation}
%   in $\Endtrl(\Prlstf)$, where the $2$-cell is the Beck--Chevalley morphism.
% \end{proof}
% Combining~\cref{lem: reidemeister trace computation} with~\cref{cor: reidemeister lefschetz triangle}, we find
% \begin{thm}
%   Let $g\colon \anima\to \anima$ be a self-map of a compact anima. There is a commutative triangle
%   \begin{equation}
%     \label{eq:reidemeister-lefschetz-triangle-with-identification}
%     \begin{tikzcd}[column sep=tiny]
%       & \SS\lbrack\cL^{g}\anima \rbrack\arrow[dr] \\
%       \SS\arrow[rr]\arrow[ru] && \SS,
%     \end{tikzcd}
%   \end{equation}
%   in which the bottom arrow is the Lefschetz trace and the arrow $\SS\lbrack\cL^{g}\anima \rbrack\to\SS$ is induced by the projection $\cL^{g}\anima \to\pt$.
% \end{thm}
% \begin{proof}
%   The second statement follows from the naturality of the identification in~\cref{lem: reidemeister trace computation}.
% \end{proof}
% We have constructed a homotopy invariant for self-maps that refines the Lefschetz trace:
% \begin{defn}
%   The map $\mdef{\rho_g}\colon \SS\to\SS\lbrack\cL^{g}\anima\rbrack$ in~\eqref{eq:reidemeister-lefschetz-triangle-with-identification} is the
%   \tdef{Reidemeister trace} of $g$.
% \end{defn}

We now apply the results of the previous section to self maps in 
topology.

\subsection{Topological fixed points}

\begin{defn}
    A \mdef{topological set} is a set $\tset$ equipped
    with a topology that is locally compact and Hausdorff. We denote
    the category of topological sets by $\Top$. 
\end{defn}

We now put ourselves in the following setting:

\begin{set}\label{set: tset}
  Consider the six-functor formalism on the geometric context
  $(\Top, \all)$ induced via the
  Liu--Zheng construction by the symmetric monoidal functor
  \[ \D{-}\colon \Top^\op \to \Prlst\,,\quad \tset \mapsto \D{\tset}=\Shv(\tset,\Sp) \]
  which sends a map $f\colon Y\to \tset$ to the pullback functor $f^*\colon\Shv (X;\Sp )\to\Shv (Y;\Sp )$. \todo{Ref Volpe?}
  By \todo{CITe } the class of all morphisms admits a suitable decomposition given by
  $(\cI ,\cP ) = (\textrm{open immersions}, \textrm{proper maps})$
  so that $\cD$ refines a to a symmetric monoidal $(\infty,2)$-functor
  \[ \cD\colon \sSpan{\Top}_{\cI,\cP}\to \Prlstf \]
  \end{set}

  Then for any $\tset \in \Top$, we have that $\tset_!\tset^* \SS= \cptSections{\tset;\SS}$
  i.e.~the spectrum $\tset_!\tset^* \SS$ is given by the
  compactly supported cohomology of the constant sheaf of spectra $\tset^*\SS$. 
  Moreover, any $\tset \in \Top$ is compact iff the map $\tset \to \pt$ is proper and any map of a compact topological sets is automatically proper. 
  In this situation, the pullback map of \cref{cons: pb tf and trc} 
  identifies with the
  classical pullback of (compactly supported) cohomology classes
  \[ \pb{f}\colon \cptSections{\tset;\SS}\to \cptSections{\tset;\SS}\,.\]
  
  Assume additionally that $\tset$ is $\cD$-suave \footnote{This holds, for instance, if $\tset$ is locally of constant shape, or more specifically if $\tset$ is an absolute neighborhood retract.}.\todo{Add REF} 
  Under these assumptions, any self map $f\colon \tset \to \tset$
  gives rise to a $\cD$-traceable span \todo{Might want to justify this}
  $\Lambda =\left(\tset= \tset \xrightarrow{f}\tset \right)$ which, as discussed
  in \cref{rmk: about pb tf and trc} has $\Sp$-linear trace simply 
  given by the symmetric monoidal trace of $\pb{f}$. 
  Moreover, \cref{prop: coincidence triangle} now supplies a commutative triangle
  \begin{equation}
    \label{eq:topological-lefschetz-trace}
    \begin{tikzcd}
      & \Gamma_c(X^f;\SS )\arrow[dr, "\lambda_f"]\\
      \SS\arrow[rr, "\tr{\pb{f}}{}"] \arrow[ru]&&\SS,
    \end{tikzcd}
  \end{equation}
  in which $\SS\to\Gamma_c(X^f;\SS )\simeq\Gamma (X^f;\SS )$ is the unit for the sheaf cohomology ring $\Gamma (X^f;\SS )\in\CAlg (\Sp)$. If $X$ is locally of constant shape, then $(\Gamma (X;\SS ),f\curvearrowright \Gamma (X;\SS ))$ is equivalent to the map induced on (singular) spherical cochains $(\SS^{\ani X},f\curvearrowright \SS^{\ani X})$. This is the usual Lefschetz trace, as one may deduce from the symmetric monoidality of the functors
  \[
    \Sp\xrightarrow{\QQ\otimes_{\SS}{-}}
    \Mod_{\QQ}\xrightarrow{\pi_*}
    \grVect_\QQ,
  \]
  where $\grVect_\QQ$ is given the Koszul rule symmetric monoidal structure. In particular, we deduce:
  
  \begin{cor}[Lefschetz fixed point theorem]
    Let $X$ be a compact Hausdorff space that is locally of constant shape, and let $f\curvearrowright X$ be some self-map. If
    \[
      \tr{\pb{f}}{} = \sum_i (-1)^i\tr{f^*\curvearrowright H^i(X;\QQ )}{}\neq 0,
    \]
    then $f$ has a fixed point.
  \end{cor}
  \begin{proof}
    If $X^f = \varnothing$ then $\Gamma_c(X^f;\SS )\simeq 0$, so commutativity of~\labelcref{eq:topological-lefschetz-trace} implies that $\tr{\pb{f}}{}\simeq 0$.
  \end{proof}
  \Onote{Maybe say something brief about fixed point indices here}

\begin{rmk}
    Suppose $X$ and $X^f$ are nice \todo{What does that mean?}. 
    Then $\lambda_f$ defines
    a class in $\pi_0\Gamma(X^f;\SS)^{\vee} \simeq H_0(X^f;\ZZ)=\ZZ[\pi_0X^f]$, i.e.~it is given by a formal sum of integers
    \[ \lambda_f= \sum_{x\in \pi_0X}[x]\lambda_f^{x}\]
    indexed by the connected components of the set of fixed points $X^f$. 
\end{rmk}
  

\subsection{Homotopy fixed points}

\begin{set}\label{set: spaces}
Consider the six-functor formalism of parametrized spectra on the $\infty$-category of spaces $\spaces$, induced via the Liu--Zheng construction by the symmetric monoidal functor
\[
    \Spc^\op\to \Prlst; \quad A \mapsto \Sp^A,
\]
which sends a map $f\colon A\to B$ to the pullback functor $f^*\colon \Sp^B \to \Sp^A $. That is, we take $(\cC ,\cE ) = (\Spc ,\all )$ and $(\cI ,\cP ) = (\textrm{all maps}, \textrm{equivalences})$ and obtain a symmetric monoidal
$(\infty,2)$-functor:
\[ \Sp^{(-)}\colon \Span{\Spc,\all, \simeq} \too \Prlstf\,.\]
\end{set}

In \cref{set: spaces} for any space 
$A\in \Spc$ we have that $A_!A^* \SS = \SS[A]=\Sigma^{\infty}_+$ is given by the spherical chains of $A$. Moreover, for any map of spaces $f\colon A \to B$, the
induced transfer map $\tf{f}\colon \SS[A]\to \SS[B]$ of
\cref{cons: pb tf and trc} is simply the pushforward map 
on the 

Now given any self map $f\colon A\to A \in \Spc$ we write
$\mdef{\cL^fA} = A^f$ for the (homotopy) fixed points of $f$, i.e.~the space of pairs $(a,\gamma )$, where $a$ is a point in $A$ and $\gamma$ is a path $a\to f(a)$. 
Now suppose that $A$ is a compact space an write $p\colon \cL^fA\to \pt$ for the terminal ap.
Then \cref{prop: coincidence triangle} supplies a commutative triangle of the form:
  \begin{equation}
    \label{eq:homotopical-lefschetz-trace}
    \begin{tikzcd}
      & \SS\lbrack \cL^fA \rbrack\arrow[dr, "\tf{p}"]\\
      \SS\arrow[rr, "\tr{\tf{f}}{}"] \arrow[ru, "\rho_f" ]&&\SS,
    \end{tikzcd}
  \end{equation}
  where the Reidemeister character $\rho_f \colon \SS\to\SS\lbrack \cL^f\anima\rbrack$ is the invariant usually referred to as the ``Reidemeister trace'' of $f$. 
  
\paragraph{Parametrized $\infty$-groupoids}
  Let $\tset$ be a locally compact Hausdorff space such that the $\infty$-topos $\Shv (\tset )$ is hypercomplete. There is a symmetric monoidal functor
  \[
    \Shv (\tset )^\op\too\etaletopoi\too\Prlst,
  \]
  where the first functor is as in~\cite[Rem~6.3.5.10]{Lurie2009} and the second functor is $\Shv ({-};\Sp )$. The composition sends a map of sheaves $f\colon G\to F$ to the pullback functor induced as in Remark 6.3.5.18 of op.~cit.~by the \'etale geometric morphism $f^*=G\times_F{-}\colon\Shv (\tset )_{/F}\to\Shv (\tset )_{/G}$. We will abuse notation by denoting the induced pullback functor by $f^*\colon\Shv (\Shv (\tset )_{/F};\Sp)\to\Shv (\Shv (\tset )_{/G};\Sp )$ also. Since this functor is induced by an \'etale geometric morphism, it has both a left adjoint $f_!$ and a right adjoint $f_*$. 
  \todo{In fact such a functor exists for any $\infty$-topos (find ref), but generally does not extend to a 6FF.}
  
  We claim that the geometric setup $(\Shv (\tset ),\textrm{all})$, the functor $\Shv (\Shv (\tset )_{/{-}};\Sp )$, and the the suitable decomposition $(\cI ,\cP ) = (\textrm{all maps},\textrm{equivalences})$ satisfy the conditions of~\cite[Prop~A.5.10]{Mann2022}. All of these are immediate, except condition (a), i.e.~\'etale base change. That is, we must show that for a commutative square
  \[
    \begin{tikzcd}
      F\arrow[r,"f"]\arrow[d,"h"] & F'\arrow[d,"h'"]\\
      G\arrow[r,"g"] & G',
    \end{tikzcd}
  \]
  the Beck--Chevalley map $f_!h^*\to (h')^*g_!$ in $\Shv (\Shv (\tset )_{/F'};\Sp )$ is an equivalence. We do this by inspecting stalks \todo{do this argument}
  
  Let $F$ be a locally constant sheaf of finitely-dominated anima on $\tset$. Then $F$ is \'etale (as all objects are) and prim with respect to the six-functor formalism described above.\todo{Prove local systems of finitely-dominated anima are prim.} Let $f\curvearrowright F$ be an endomorphism. We then get a commutative triangle
  \begin{equation}
    \label{eq:homotopical-lefschetz-trace}
    \begin{tikzcd}
      & \SS\lbrack L^fF\rbrack\arrow[dr]\\
      \SS_\tset\arrow[rr,"\chr {}f"] \arrow[ru]&&\SS_\tset,
    \end{tikzcd}
  \end{equation}
  \todo{work out what this looks like on stalks.}

 
% \newpage

% Let $\Spc$ denote the $\infty$-category of spaces. For any $X \in \Spc$, we write
% $\Sp^{X}$ for the category of local systems of spectra on $X$. For any map of spaces
% $f\colon X \to Y$, the pullback functor $f^{\ast}\colon \Sp^Y\to \Sp^X$ admits both adjoints
% $f_{\sharp}\dashv f^{\ast} \dashv f_{\ast}$, given by left and right Kan extension respectively.
% As before, if $Y=\pt$ and $f$ is the terminal map, we also write $X^{\ast}=f^{\ast}$ and so on.
% With this notation, we have that the composite
% \[X_{\sharp}X^{\ast}\SS = \SS[X] = \colim_X \SS\]
% is given by the spherical chains on $X\in \Spc$.
% By \todo{CITE}, the assignment taking
% a span of spaces $X\xleftarrow{f}Y \xrightarrow{g} Z$ to the composite functor
% $\Sp^X \xrightarrow{f^{\ast}} \Sp^Y \xrightarrow{g_{\sharp}} \Sp^Z$ refines to a symmetric monoidal
% functor
% \[ (\Sp^{(-)})^{\ast}_{\sharp}\colon \Span{\Spc} \too \Prlst\,.\]
% Given a self map $f\colon X\to X$ of a space $X\in \Spc$, we write $\cL^{f}X$ for the pullback
% \[\begin{tikzcd}
%     \cL^{f}X  & X \\
%     X & X\times X\,.
%     \arrow[from=1-1, to=1-2]
%     \arrow[from=1-1, to=2-1]
%     \arrow["\Delta", from=1-2, to=2-2]
%     \arrow["\Gamma^f"', from=2-1, to=2-2]
%   \end{tikzcd}\]
% Note that, $\cL X = \cL^{\id}X= \Map(S^1,X)$ is simply the free loop space of $X$. Arguing as in
% the proof of \cref{lem: trace computation}, we immediately learn that
% \[ \tr{f^{\ast}\colon \Sp^X\to \Sp^X}{\L}\simeq \SS[\cL^{f}X] \in \Sp\,.\]

% In what follows, fix a self map $f\colon X\to X$ of a compact space $X\in \Spc^{\omega}$

% \begin{prop}\label{prop: weak lefschetz homotopical}
%   The trace $\tr{\SS[f]}{}\colon \SS\to \SS$ factors through the pushforward
%   $\SS[\cL^{f}X] \to \SS$ induced by the terminal map $\cL^{f}X\to \pt$.
% \end{prop}
% \begin{proof}
%   Since $X$ is compact, we have a chain of adjoints
%   $X_{\sharp} \dashv X^{\ast} \dashv X_{\ast} \dashv X^!$, so that $X_{\sharp}$ and $X^{\ast}$
%   are both strongly continuous. Let $\tau\colon X_{\sharp} f^{\ast}\to X_{\sharp}$ be the natural
%   transformation obtained by whiskering the counit $f_{\sharp}f^{\ast} \to \id_{\Sp^X}$ with
%   $X_{\sharp}$. Moreover, writing $\can \colon X^{\ast}\xrightarrow{\sim}f^{\ast}X^{\ast}$ for
%   the natural equivalence, we obtain a diagram in $\Endtrl(\Prlstf)$ of the form:
%   \[\begin{tikzcd}
%       {\Sp} & {\Sp^X} & {\Sp} \\
%       {\Sp} & {\Sp^X} & {\Sp}
%       \arrow["{X^{\ast}}", from=1-1, to=1-2]
%       \arrow[equals, from=1-1, to=2-1]
%       \arrow["{X_{\sharp}}", from=1-2, to=1-3]
%       \arrow["f^{\ast}",from=1-2, to=2-2]
%       \arrow[equals, from=1-3, to=2-3]
%       \arrow["\mathrm{can}"', Rightarrow, from=2-1, to=1-2]
%       \arrow["{X^\ast}"', from=2-1, to=2-2]
%       \arrow["{\tau}"', Rightarrow, from=2-2, to=1-3]
%       \arrow["{X_{\sharp}}"', from=2-2, to=2-3]\,.
%     \end{tikzcd}\]
%   Unwinding the definitions, we see that the composite is a map $(\Sp,\id)\to (\Sp,\id)$ given
%   by tensoring with the spectrum $\SS[X]$, together with the natural endomorphism of
%   $-\otimes \SS[X]$ given by the pushforward map $\SS[f]\colon \SS[X] \to \SS[X]$. Thus,
%   applying the functor $\TrcL$ to the diagram above, we obtain a commutative diagram
%   in spectra
%   \[\begin{tikzcd}
%       & {\tr{f^{\ast}}{\L}} \\
%       \SS && \SS
%       \arrow["{\TrcL(X_{\sharp},\tau)}", from=1-2, to=2-3]
%       \arrow["{\TrcL(X^{\ast},\can)}", from=2-1, to=1-2]
%       \arrow["{\tr{\SS[f]}{}}"', from=2-1, to=2-3]\,,
%     \end{tikzcd}\]
%   so it suffices to show that, under the identification $\tr{f^{\ast}}{\L} \simeq \SS[\cL^{f}X]$
%   the map $\TrcL(X_{\sharp},\tau)$ is given by the pushforward along the terminal map,
%   which again follows exactly as in \cref{lem: trace computation}
% \end{proof}

% \begin{defn}
%   Given a self map of a compact space $f\colon X\to X$, 
%   we denote the homology class factoring the trace obtained in
%   \cref{prop: weak lefschetz homotopical} by
%   \[ \mdef{\rho_f} \colon \SS\to \SS[\cL^{f}X]\,,\]
%   and call $\rho_f$ the \tdef{Reidemeister trace} of $f$.
% \end{defn}

% \subsection{Cosheaves and assembly}

% For a locally compact Hausdorff space $\tset \in \Top$, we write
% \[ \ani{\tset}\coloneq \ani{\mathrm{Sing}_{\bullet}{\tset}} \in \Spc\]
% for the geometric realization of the singular complex of $\tset$.
% If $\tset$ is compact and suave, we have a natural identification
% \[ \Gamma(\tset;\omega_\tset) \simeq \SS[\ani{\tset}]\]
% of the global sections of the dualizing sheaf with the spherical chains on $\ani{\tset}$,
% which we thus simply denote $\SS[\tset]$. 

% We now relate the Lefschetz index of a map $f\colon \tset \to \tset$ to the Reidemeister
% trace of the induced map $\ani{f}\colon \ani{\tset} \to \ani{\tset}$. Denote by
% $i\colon \ani{A^f} \to \cL^{\ani{f}}\ani{A}$ the natural map induced by the assignment
% taking a fixed point $a\in A^f$ to the pair $(a,\mathrm{const}_a)$ were $\mathrm{const}_a$
% is the constant path $a\to f(a)=a$. 

% \begin{thm}
%   For any self map of a compact suave Hausdorff space $f\colon \tset \to \tset$,
%   the Reidemeister trace $\rho_{\ani{f}}$ factors naturally in $f$ through the Lefschetz trace,
%   i.e.~we have a commutative diagram in spectra of the form:
% \[\begin{tikzcd}
% 	& {\Sections{\tset^f;\omega_{\tset^f}}} \\
% 	\SS & {\SS[\cL^{f}\ani{\tset}]}
% 	\arrow[from=1-2, to=2-2]
% 	\arrow["{\ch{}{f}}", from=2-1, to=1-2]
% 	\arrow["{\rho_{\ani{f}}}"', from=2-1, to=2-2]
% \end{tikzcd}\]
% \end{thm}

% We prove this by constructing a factorization of the map
% $(\ani{\tset}^{\ast},\can)\colon (\Sp,\id)\to (\Sp^{\ani{\tset}},\abs{f}^{\ast})$ in
% the category $\Endtrl(\Prlstf)$.

% \begin{defn}
%   Denote by $\Prdualf$ the sub-$(\infty,2)$-category of $\Prlstf$ spanned by the dualizable objects
%   in $\Prlstf$ with strongly continuous functors between them.
% For any $\tset \in \Top$ we write
% \[ \coShat(\tset) \coloneq \map_{\Prdualf}(\D{\tset}, \Sp)\]
% for the weak dual of $\D{\tset}= \Shv(\tset; \Sp)$ in $\Prdualf$.
% \end{defn}
% Barthels--Nikolaus construct a factorization of the functor $\ani{\tset}^{\ast}\colon \Sp \to \Sp^{\ani{\tset}}$
% of the form
% \[ \Sp \xrightarrow{\chi_{\tset}^{\loc}} \coShat(\tset) \xrightarrow{\ass_{\tset}} \Sp^{\ani{\tset}}\,,\]
% where $\chi_{\tset}^{\loc}(\SS) $ is cosheaf taking an open $U\subseteq \tset$ to the spectrum
% $\SS[\abs{U}]$.

% \begin{prop}
% We have a natural factorization in $\Endtrl(\Prlstf)$ of the form:
% \[\begin{tikzcd}
% 	\Sp & {\coShat{\tset}} & {\Sp^{\ani{\tset}}} \\
% 	\Sp & {\coShat{\tset}} & {\Sp^{\ani{\tset}}}
% 	\arrow["{\chi_{\tset}^{\loc}}", from=1-1, to=1-2]
% 	\arrow["\id", from=1-1, to=2-1]
% 	\arrow["\ass_{\tset}", from=1-2, to=1-3]
% 	\arrow["{f^{\natural}}", from=1-2, to=2-2]
% 	\arrow["{f^{\ast}}", from=1-3, to=2-3]
% 	\arrow[ Rightarrow, from=2-1, to=1-2]
% 	\arrow["{\chi_{\tset}^{\loc}}"', from=2-1, to=2-2]
% 	\arrow[ Rightarrow, from=2-2, to=1-3]
% 	\arrow["\ass_{\tset}",from=2-2, to=2-3]
% \end{tikzcd}\]
% Moreover, the induced map on $\TrcL$ factors through $\TrcL(\D{\tset},f^{\ast})\simeq \Gamma_c(\tset;\SS_A)^{\vee}$. 
% \end{prop}



%%% Local Variables:
%%% mode: LaTeX
%%% TeX-master: "main"
%%% End:
